% =============================================================================
% cap3_pipeline/04_fase3_classificacao.tex
% Seção 3.4 -- Fase III: Classificação multi-espécie
%   Modelo: SpeciesNet (pré-treinado, as-is).
%   Função: discriminar a espécie do animal detectado na Fase II, filtrando
%   o "ruído biológico" (gambás, quatis, lagartos, aves, gato-mourisco) e
%   encaminhando apenas Felis catus para a re-ID da Fase IV.
%   Texto corrido, impessoal, sem listas no corpo.
% =============================================================================
\section{Fase III --- Classificação de espécie}
\label{sec:pipeline-fase3}

A Fase~III recebe da Fase~II quadros com detecções da classe macroscópica ``animal'' e tem como função única discriminar a espécie de cada detecção. Por estarem dispostos no solo e em formato aberto, os potes de ração do Campus~2 atraem rotineiramente fauna não-alvo --- gambás (\estr{Didelphis albiventris}), quatis (\estr{Nasua nasua}), lagartos teiú, pombos, aves nativas e, criticamente, o felino silvestre gato-mourisco (\estr{Herpailurus yagouaroundi})~\cite{goncalves2025wildcat,cove2022capital}. Submeter indiscriminadamente essas espécies à reidentificação individual da Fase~IV, treinada e validada apenas para \estr{Felis catus}, produziria identificações falsas e contaminação persistente do histórico da colônia.

\subsection*{Escolha do modelo: \modeloSN~como classificador base}

O modelo adotado nesta fase é o \modeloSN, classificador hierárquico de espécies treinado em centenas de datasets de \emph{camera trapping} ao redor do mundo, liberado em código aberto pelo \emph{Google Research}~\cite{google_speciesnet_blog}. Três fatores justificam sua escolha. Primeiro, foi desenhado para a tarefa exata desta fase: classificação de espécies sobre recortes de animais previamente detectados, em condições típicas de \emph{camera trapping} (iluminação variável, ângulos não controlados, ocasional NIR), e não em fotografia controlada. Segundo, sua hierarquia taxonômica interna permite consumir a saída em diferentes níveis de generalidade --- da família ao indivíduo --- o que é útil ao filtro deste pipeline: quando o modelo não chega à espécie com confiança suficiente, ainda assim entrega a família ou a ordem. Terceiro, sua disponibilidade pública e integração com a pilha \emph{PyTorch-Wildlife} usada na Fase~II tornam a adoção operacionalmente simples.

Em linha com a decisão de escopo da Seção~\ref{sec:pipeline-visao-geral}, o \modeloSN~é utilizado sem \emph{fine-tuning} sobre dados do Campus~2. A cobertura do modelo abrange com folga as espécies de fauna não-alvo observadas em campo --- gambás e quatis são classes bem representadas em datasets neotropicais incorporados ao treinamento --- e tampouco há motivo para esperar falha sistemática no reconhecimento de \estr{Felis catus}, classe abundante no treinamento global. Eventual \emph{fine-tuning} leve sobre um conjunto representativo do Campus~2 fica registrado como trabalho futuro no Capítulo~\ref{cap:conclusao}.

\subsection*{Operação da fase no pipeline}

Para cada detecção recebida da Fase~II, o \modeloSN~é aplicado sobre o recorte da \emph{bounding box} e produz uma distribuição de probabilidade sobre seu vocabulário taxonômico. A política de decisão é conservadora: apenas detecções classificadas como \estr{Felis catus} com confiança acima de um limiar configurável seguem para a Fase~IV. Quadros classificados com confiança em outra espécie são registrados em uma base de fauna não-alvo --- subproduto científico do projeto, retomado nos trabalhos futuros do Capítulo~\ref{cap:conclusao} --- mas não avançam no pipeline. Quadros sem classificação confiante (alta entropia distribuída entre várias classes plausíveis, ou baixa confiança concentrada em \estr{Felis catus}) são roteados para uma fila de validação humana, em que voluntários da AEX Gatosdoc2 podem confirmar ou rejeitar a classificação por uma interface mínima de revisão.

A fila de validação humana cumpre dupla função: impede que casos duvidosos sigam para a re-ID e contaminem o histórico da colônia (papel operacional) e, ao longo do tempo, constitui um pequeno dataset rotulado do domínio do Campus~2 (papel metodológico), que poderia, em trabalho futuro, alimentar \emph{fine-tuning} leve do \modeloSN~ou reavaliação dos limiares de confiança das Fases~II e III. No escopo deste projeto, a fila é tratada apenas como camada de contenção, sem que as decisões algorítmicas das fases posteriores dependam dela.

\subsection*{Tratamento de espécies próximas e do gato-mourisco}

Caso particular merecedor de menção explícita é o do \estr{Herpailurus yagouaroundi}: felino silvestre nativo, de porte similar ao do gato doméstico em alguns indivíduos, presente em fragmentos urbanos da região central de São Paulo~\cite{goncalves2025wildcat}. Trata-se exatamente do cenário em que uma classificação binária ``gato versus não-gato'' seria semanticamente correta mas operacionalmente desastrosa, por misturar duas espécies distintas no estágio de re-ID. O \modeloSN, treinado sobre vocabulário taxonômico abrangente, contempla a espécie no seu repertório e é capaz de distinguir as duas em condições adequadas de iluminação e ângulo. Em condições degradadas (NIR, ângulo oblíquo, oclusão parcial), é razoável esperar saídas de baixa confiança e, portanto, encaminhamento para a fila de validação humana --- comportamento aceitável para fins do projeto.

Ao final da Fase~III, segue para a Fase~IV um subconjunto ainda menor de quadros: aqueles em que houve detecção válida de animal na Fase~II, a classificação multi-espécie convergiu de forma confiante para \estr{Felis catus} e, quando acionada, a validação humana confirmou essa classificação. É sobre esse subconjunto, e somente sobre ele, que a tarefa-fim do sistema pode operar com taxas de falso positivo aceitáveis.
